\chapter{Pascal}\label{cha:pascal}

The machine has two Pascals, and they are different animals. Turbo
Pascal~3 runs on the Tube under CP/M (Chapter~\ref{cha:cpm}): an
on-machine compiler, a Z80 program, its own world. \emph{Mad Pascal}
is the other kind --- a cross-compiler, like cc65 the ROM is built with:
you write on the desktop, \texttt{mp} compiles to 6502 assembly, MADS
assembles it, and the result is a \texttt{.prg} the shell's \verb|RUN|
loads like any other program. It is Tomasz Biela's compiler (MIT), a
Turbo-Pascal-flavoured language with 8-, 16- and 32-bit integers, fixed
and floating point, strings, records, pointers, units and inline
assembly, made for the Atari and since taught the C64, the X16 and the
Neo6502. The K4510 is its newest target.

\begin{type}
PMANDEL
\end{type}

That is \texttt{demo/pas/pmandel.pas}: the Mandelbrot set, 78 by 28
cells in the machine's sixteen colours, in a little over four seconds
at 40.5~MHz --- fixed point, 32-bit integers, forty lines of Pascal.
\verb|PSIEVE| is the classic sieve, ten passes timed by the frame
counter; \verb|HELLO| prints, shows the colours, and echoes what you
typed after its name.

\section{The target}
Everything a Pascal program needs from the machine is in
\texttt{pascal/} of the repository, laid out as it lives inside a
Mad-Pascal checkout: the runtime base (\texttt{base/rtl6502\_k4510.asm}
and \texttt{base/k4510/}), where \texttt{@putchar} writes to JIM, the
terminal (Chapter~\ref{cha:tube}); the SYSTEM and CRT units' machine
halves (\texttt{lib/*\_k4510.inc}); and a \texttt{k4510} unit. The
consequences for the programmer:

\begin{itemize}
\item \verb|Write| and \verb|WriteLn| go through JIM, so the CRT unit is
  the real thing: \verb|GotoXY|, \verb|TextColor| and
  \verb|TextBackground| (the palette's constants: \verb|BLUE|,
  \verb|YELLOW|, \verb|LIGHT_GREEN|\ldots), \verb|ClrScr|, \verb|ClrEol|,
  \verb|InsLine|/\verb|DelLine|, \verb|WhereX|/\verb|WhereY|,
  \verb|CursorOn|/\verb|CursorOff|, \verb|ReadKey| and \verb|KeyPressed|
  on the keyboard device, \verb|Delay| and \verb|Pause| on the frame
  counter, \verb|Sound| on SID~0. \verb|TextMode(0)| puts the ROM's
  screen back.
\item \verb|ParamCount| and \verb|ParamStr| come from the shell line
  (\verb|RUN HELLO one two|, or just \verb|HELLO one two|);
  \verb|ParamStr(0)| is the whole tail.
\item \verb|uses k4510| gives every chip as a typed variable at its
  address --- \verb|VICKE[reg]|, \verb|SIDREG[]|, \verb|DMA_SRC|,
  \verb|FS_CMD|, \verb|SYS_FRAMES|, \verb|MATH_F[]|, \verb|TERM_CX|
  and the rest --- plus \verb|FarPeek|/\verb|FarPoke| into the 256~MB
  through the 45GS02's flat addressing, \verb|DmaCopy|/\verb|DmaFill|,
  \verb|Shell('DIR')|, \verb|LoadFile|/\verb|SaveFile| through the ROM,
  \verb|WaitVBlank|, and \verb|TermWrite| for raw escape sequences.
\item Memory: code from \texttt{\$0800}, the program's own; zero page
  \texttt{\$22--\$3F} for the compiler's registers and
  \texttt{\$64--\$A3} for its expression stack (the ROM keeps
  \texttt{\$02--\$21} and \texttt{\$F0--\$F9}); \texttt{\$0300} the string
  buffer. Programs return to the shell with the ROM's state intact.
\end{itemize}

\section{Building}
On the desktop, with FPC and a Mad-Pascal checkout (the installer
patches its target table and rebuilds \texttt{mp}):

\begin{verbatim}
make pascal-install      # once: MP_DIR=... (default ~/Projects/neo6502_dev/Mad-Pascal)
make pascal              # demo/pas/*.pas -> fs/PRG/*.prg
\end{verbatim}

A new program is a file in \texttt{demo/pas/}; \texttt{make} finds it.
The compiled \texttt{.prg} files are committed, so a machine without
the toolchain still runs them, and the test suite (\texttt{test/pastest.sh})
runs \verb|HELLO| and \verb|PSIEVE| through the ROM.

\begin{aside}
\textbf{Not yet:} floating point still runs in software (Mad Pascal's own
IEEE single). The MATH unit at \texttt{\$D700} does the same arithmetic
in one register write; wiring the compiler's \texttt{single} to it is
the next step on this road, and the graph unit on VICKe's bitmap layer
the one after.
\end{aside}
